\begin{frame}{PageRank의 핵심 아이디어와 정식화}
  \begin{itemize}
    \item 이 그래프 위를 \textbf{무작위로 영원히 돌아다닌다면},
      각 노드에 머물러 있을 확률은 각각 얼마나 될까?
    \item 그 답 = 그 노드의 ``\textbf{중요도}'' (1998, Brin \& Page).
    \item 많은 페이지가 링크하는 페이지일수록, 또 \textbf{중요한
      페이지가} 링크하는 페이지일수록, \kb{무작위 서퍼}(random
      surfer)가 더 자주 머문다.
  \end{itemize}
  \vskip0.5em
  \begin{itemize}
    \item ``누가 누구를 링크하는가''를 \textbf{무작위 서퍼의 이동
      행렬}로 정식화 $\to$ 1998년 논문
      \textit{``The Anatomy of a Large-Scale Hypertextual Web Search
      Engine''}이 실용 알고리즘으로 완성.
  \end{itemize}
  {\footnotesize (사명 ``Google''은 $10^{100}$을 뜻하는
  ``googol''에서 온 말장난 --- 조직하려던 웹의 규모에 대한 선언.)}
\end{frame}
